\documentclass[10pt,a4paper]{amsart}
\usepackage[margin=19mm]{geometry}
\usepackage{amsmath,amssymb,mathtools}
\usepackage[scr=rsfs,cal=euler]{mathalfa}
\usepackage{xfrac}
\usepackage[unicode,hidelinks,bookmarks=false]{hyperref}
\newcommand{\ie}{i.e.\ }
\newcommand{\cf}{cf.\ }
\newcommand{\resp}{respectively\ }
\newcommand{\Subsection}{\subsection}
\providecommand{\nobreakdash}{\nobreak\hbox{-}}
\setlength{\emergencystretch}{2em}
\multlinegap0pt
\usepackage{bm}
\usepackage{bbm}
\usepackage{dsfont}
\usepackage{xspace}
\usepackage{cancel}
\usepackage{mathtools}
\usepackage{amsthm}
\makeatletter
\@ifundefined{thm}{\newtheorem{thm}{Thm}}{}
\@ifundefined{prop}{\newtheorem{prop}[thm]{Proposition}}{}
\makeatother
\title[Minimal Surface Entropy and Average Area Ratio]{Minimal Surface Entropy and Average Area Ratio}
\author{Ben Lowe, Andre Neves}
\date{Source: arXiv:2110.09451v3 (CC BY 4.0)}
\begin{document}
\maketitle

\noindent\emph{Source and licence.} Minimal Surface Entropy and Average Area Ratio. Version arXiv:2110.09451v3 (CC BY 4.0), distributed under the Creative Commons Attribution 4.0 International licence, as recorded in the arXiv licence metadata for this exact version and shown on the abstract page https://arxiv.org/abs/2110.09451v3. Full rights notice: licenses/arxiv-2110.09451-CC-BY-4.0.txt.\par\smallskip

\noindent\textit{Editorial scope.} Counted unit: the Prop whose source label is \texttt{asymptotic.expn} in the article (source0.tex lines 966--972, source label \texttt{asymptotic.expn}), delivered together with the article's own complete original proof of it (source0.tex lines 974--1031, one \texttt{proof} environment of 58 lines). No mathematics is written, repaired or re-derived: statement and proof are the article's own text line for line. Required context retained uncounted: none. Every definition, display and label of those lines is retained unchanged. Omitted from this package and not redistributed: 1--965, 973--973, 1032--1249. Nothing inside a retained excerpt is omitted or altered: every retained body is delivered exactly as the article's lines are written, so no omission receipt of any kind accompanies this package. Citations: the retained excerpts carry no citation command at all, so no bibliography entry of the article is transcribed and no reference is redistributed. Reference adaptation: none was needed and none was made; every reference inside the delivered unit resolves inside it, so no unresolved reference or citation is produced. Printed slips kept literal: no printed word, symbol, hypothesis or number of the counted statement or of its proof was altered. Layout and class: the article's own class is \texttt{amsart} with packages including verbatim, latexsym, amssymb, amsmath, color, epsfig, hyperref; the wrapper is the lane's pinned \texttt{amsart} editorial wrapper at 10pt on A4, which restates the theorem-like environments the retained text needs and adds the article's own macro definitions that the retained text needs (none), with extra wrapper packages (bm, bbm, dsfont, xspace, cancel, mathtools). The delivered numbering is the unit's own. Assets: no retained span contains a figure environment, a graphics-inclusion command, a PDF-page inclusion, a table or a TikZ picture; no image file is copied into this package, the package declares no asset dependency and no image file is redistributed with it. Licence: version arXiv:2110.09451v3 of this article is distributed under the Creative Commons Attribution 4.0 International licence, as recorded in the arXiv licence metadata for this exact version and shown on the abstract page https://arxiv.org/abs/2110.09451v3. A full rights notice is delivered at licenses/arxiv-2110.09451-CC-BY-4.0.txt. Characters: every retained excerpt is pure ASCII, so the delivered engine, XeLaTeX with its default Latin Modern text font, reports no missing character.
	\begin{prop}\label{asymptotic.expn}
		With $\gamma\in \mathcal C_{\varepsilon}$, $x\in D(\gamma)$ and $y=F_{\gamma}(x)$ we have
		\begin{multline*}Ric(g)_y(\nu_g(\gamma),\nu_g(\gamma))-\frac{1}{3}R(g)(y)=L(h)_x(n(\gamma),n(\gamma))\\ +O^3_x(|h|^2)+O^1_x(|f_{\gamma}|^2).
		\end{multline*}
		With $|Jac_g F_{\gamma}|$ the Jacobian of $F_{\gamma}:(D(\gamma),g_0)\rightarrow (\Sigma_{g}(\gamma),g)$ we have
		$$|Jac_g F_{\gamma}|(x)=1+O_x^1(|h|)+O_x^1(|f_{\gamma}|).$$
	\end{prop}

	\begin{proof}
		Set $$X:=\{(\gamma,f):\gamma\in\mathcal C_{\bar \varepsilon}, f\in C^2(D(\gamma)), |f|_{C^2}\leq 1\}.$$
		Given $(\gamma,f)\in X$ set
		\begin{equation*}
			F(\gamma,f): D(\gamma)\rightarrow \Sigma_{g}(\gamma),\quad x\mapsto \cosh(f(x))x+\sinh(f(x))n(\gamma)(x)
		\end{equation*}
		and $\Sigma(\gamma,f):=F(\gamma,f)(D(\gamma))$. Let $n(\gamma,f)$ be the unit normal vector field with respect to $g_0$ along $\Sigma(\gamma,f)$, defined so that $n(\gamma,f)$ depends smoothly on $(\gamma,f)$ and $n(\gamma,0)=n(\gamma)$.
		We have for all $0\leq t\leq 1$
		$$\left|\frac{d^k}{(dt)^k}(n(\gamma,tf)(F(\gamma,tf)(x))\right|=O_x^1(|f|^k), \quad k=1,2.$$
		Hence, if $G$ is a $2$-tensor on $M$ with $|G|_{C^2}\leq 1$, then
		$$\alpha(t):=G(n(\gamma,tf)(F(\gamma,tf)(x)),n(\gamma,tf)(F(\gamma,tf)(x)))$$
		satisfies 
		$$|\alpha'(0)|\leq (|\nabla G|(x)+|G|(x))O_x^1(|f|)\quad\mbox{and}\quad \sup_{0\leq t\leq 1}|\alpha''(t)|=O_x^1(|f|^2).$$
		Therefore, with $y=F(\gamma,f)(x)$, we obtain from Taylor's expansion
		\begin{multline}\label{taylor.epn}
			|G(n(\gamma,f)(y),n(\gamma,f)(y))-G(n(\gamma)(x),n(\gamma)(x))|\\
			\leq |\nabla G|^2(x)+|G(x)|^2+O_x^1(|f|^2).
		\end{multline}
		Setting $G=|\nabla^k h|^2g_0$ in this inequality we deduce that for  $k=0,1,2.$
		\begin{equation}\label{base.point}
			|\nabla^k h|(y)=O^{k+1}_x(|h|)+O_x^1(|f|).
		\end{equation}
		Setting  $G=L(h)$ in \eqref{taylor.epn} and using the fact that $L$ is a second order differential operator we deduce
		\begin{multline}\label{expn.Lh}
			|L(h)_x(n(\gamma)(x),n(\gamma)(x))-L(h)_y(n(\gamma,f)(y),n(\gamma,f)(y))|\\
			=O_x^3(|h|^2)+ O_x^1(|f|^2).
		\end{multline}
		Let $\nu_g(\gamma,f)$ denote the unit normal vector field along $\Sigma(\gamma,f)$ with respect to $g=g_0+h$ so that $\nu_{g_0}(\gamma,f)=n(\gamma,f)$ and $\nu_g$ depends smoothly on its parameters. Then
		$$|\nu_g(\gamma,f)(y)-n(\gamma,f)(y)|=O_y^0(|h|)=O_x^1(|h|)+ O_x^1(|f|),$$
		where in the last identity we used \eqref{base.point}. Using this identity and \eqref{base.point}  we have
		\begin{multline*}
			L(h)_y(\nu_g(\gamma,f)(y),\nu_g(\gamma,f)(y))\\
			=L(h)_y(n(\gamma,f)(y),n(\gamma,f)(y))+O_x^3(|h|^2)+ O_x^1(|f|^2).
		\end{multline*}
		Combining with \eqref{expn.Lh} we deduce
		\begin{multline*}
			L(h)_y(\nu_g(\gamma,f)(y),\nu_g(\gamma,f)(y))\\
			=L(h)_x(n(\gamma)(x),n(\gamma)(x))+O_x^3(|h|^2)+ O_x^1(|f|^2).
		\end{multline*}
		Using Taylor's expansion  we have  that for every $y\in M$, every unit vector field $Y\in T_yM$ and $g_0+h\in \mathcal U$
		$$|\mathring{Ric}(g_0+h)_y(Y,Y)-L(h)_y(Y,Y)|=O_y^2(|h|^2)=O_x^3(|h|^2)+O_x^1(|f|^2).$$
		Therefore
		\begin{multline*}
			\mathring{Ric}(g_0+h)_y(\nu_g(\gamma,f)(y),\nu_g(\gamma,f)(y))\\
			=L(h)_y(n(\gamma)(x),n(\gamma)(x))+O_x^3(|h|^2)+ O_x^1(|f|^2).
		\end{multline*}
		The first statement in the proposition follows from choosing $f=f_\gamma$ in the identity above.
		
		We now prove the statement regarding $|Jac_g F_{\gamma}|$. With $(\gamma,f)\in X$,  denote by $|Jac_g F(\gamma,f)|(g)$ the Jacobian of $F(\gamma,f):(D(\gamma),g_0)\rightarrow (\Sigma(\gamma,f),g)$. With $g=g_0+h\in \mathcal U$ and $y=F(\gamma,f)(x)$ we have
		\begin{align*}
			|Jac_g F(\gamma,f)|(x) & =|Jac_{g_0}F(\gamma,f)|(x)+ O^0_y(|h|)\\
			&=|Jac_{g_0}F(\gamma,f)|(x)+ O^1_x(|h|)+O^1_x(|f|)\\
			&=|Jac_{g_0}F(\gamma,0)|(x)+O^1_x(|h|)+O^1_x(|f|)\\
			&=1+O^1_x(|h|)+O^1_x(|f|).
		\end{align*}
		
		
	\end{proof}

\end{document}
